\documentclass[11pt,letterpaper]{article}
\usepackage[T1]{fontenc}
\usepackage{lmodern,microtype}
\usepackage[margin=1in]{geometry}
\usepackage[pdfauthor={Christopher Mount},pdftitle={Teaching Philosophy}]{hyperref}
\setlength{\parindent}{0pt}
\setlength{\parskip}{.5\baselineskip}
\raggedbottom

\begin{document}
\begin{center}
  {\Large\bfseries Teaching Philosophy}\par
  Christopher Mount
\end{center}

My responsibility as a mathematics instructor is to help students decide how to approach a problem, explain their choices, and assess the result. These abilities matter when students encounter unfamiliar applications, consult reference material, or use computational tools. They also guide how I think about teaching in the presence of AI. Access to a solution makes it increasingly important to establish what a student understands about the reasoning that produced it.

I distinguish reproducing a method from adapting it. A student who has learned the integral test should be able to apply it to familiar examples. A deeper understanding appears when the associated integral resists direct evaluation and the student recognizes that comparison can still establish convergence. I want students to experience this progression repeatedly, using established ideas to learn a new method and later using that method to understand something further.

Clear instruction and opportunities to work through difficulties both contribute to this process. I value lectures that develop an argument with a clear direction, followed by time for students to attempt problems with guidance. In a course I design, I would favor a weekly rhythm that includes sustained lectures and a dedicated session for facilitated group work. This would give students time to see how ideas fit together and then test their understanding through practice and discussion.

My role during that practice is often to ask a question that reconnects students with something they know. When a student hesitates over a trigonometric identity, I may ask a tablemate to identify the relevant property of the function. When a student struggles with a Taylor expansion, I may direct them back to a familiar series and ask what substitution would accomplish. I want students to participate in constructing the explanation. I have had to work on my own patience here, allowing several seconds for an answer before stepping in.

I also try to make uncertainty an ordinary part of mathematical work. I acknowledge when I need to check something myself. Students should be able to ask questions, consult their notes, and revise an approach without treating those actions as evidence that they cannot learn mathematics. My grading similarly aims to identify what students understand, giving the greatest weight to the central concept a problem assesses.

AI requires us to reconsider what assignments reveal about learning. I expect students preparing for STEM careers to need experience using it, including practice checking its claims and recognizing its limitations. At the same time, in-person quizzes, exams, and mathematical conversations provide evidence of what students can explain and do independently. I am interested in how brief oral explanations might contribute to assessment, with careful attention to the anxiety they can create.

My own use of AI includes interactive teaching materials, simulations, and assistance organizing a database of examination problems. I write my exams and have used AI to check them for typographical errors. This experience informs my view that learning to use a tool includes learning which decisions remain the user's responsibility.

\end{document}
